% Options for packages loaded elsewhere
\PassOptionsToPackage{unicode}{hyperref}
\PassOptionsToPackage{hyphens}{url}
\documentclass[
]{article}
\usepackage{xcolor}
\usepackage[margin=1in]{geometry}
\usepackage{amsmath,amssymb}
\setcounter{secnumdepth}{-\maxdimen} % remove section numbering
\usepackage{iftex}
\ifPDFTeX
  \usepackage[T1]{fontenc}
  \usepackage[utf8]{inputenc}
  \usepackage{textcomp} % provide euro and other symbols
\else % if luatex or xetex
  \usepackage{unicode-math} % this also loads fontspec
  \defaultfontfeatures{Scale=MatchLowercase}
  \defaultfontfeatures[\rmfamily]{Ligatures=TeX,Scale=1}
\fi
\usepackage{lmodern}
\ifPDFTeX\else
  % xetex/luatex font selection
\fi
% Use upquote if available, for straight quotes in verbatim environments
\IfFileExists{upquote.sty}{\usepackage{upquote}}{}
\IfFileExists{microtype.sty}{% use microtype if available
  \usepackage[]{microtype}
  \UseMicrotypeSet[protrusion]{basicmath} % disable protrusion for tt fonts
}{}
\makeatletter
\@ifundefined{KOMAClassName}{% if non-KOMA class
  \IfFileExists{parskip.sty}{%
    \usepackage{parskip}
  }{% else
    \setlength{\parindent}{0pt}
    \setlength{\parskip}{6pt plus 2pt minus 1pt}}
}{% if KOMA class
  \KOMAoptions{parskip=half}}
\makeatother
\usepackage{color}
\usepackage{fancyvrb}
\newcommand{\VerbBar}{|}
\newcommand{\VERB}{\Verb[commandchars=\\\{\}]}
\DefineVerbatimEnvironment{Highlighting}{Verbatim}{commandchars=\\\{\}}
% Add ',fontsize=\small' for more characters per line
\usepackage{framed}
\definecolor{shadecolor}{RGB}{248,248,248}
\newenvironment{Shaded}{\begin{snugshade}}{\end{snugshade}}
\newcommand{\AlertTok}[1]{\textcolor[rgb]{0.94,0.16,0.16}{#1}}
\newcommand{\AnnotationTok}[1]{\textcolor[rgb]{0.56,0.35,0.01}{\textbf{\textit{#1}}}}
\newcommand{\AttributeTok}[1]{\textcolor[rgb]{0.13,0.29,0.53}{#1}}
\newcommand{\BaseNTok}[1]{\textcolor[rgb]{0.00,0.00,0.81}{#1}}
\newcommand{\BuiltInTok}[1]{#1}
\newcommand{\CharTok}[1]{\textcolor[rgb]{0.31,0.60,0.02}{#1}}
\newcommand{\CommentTok}[1]{\textcolor[rgb]{0.56,0.35,0.01}{\textit{#1}}}
\newcommand{\CommentVarTok}[1]{\textcolor[rgb]{0.56,0.35,0.01}{\textbf{\textit{#1}}}}
\newcommand{\ConstantTok}[1]{\textcolor[rgb]{0.56,0.35,0.01}{#1}}
\newcommand{\ControlFlowTok}[1]{\textcolor[rgb]{0.13,0.29,0.53}{\textbf{#1}}}
\newcommand{\DataTypeTok}[1]{\textcolor[rgb]{0.13,0.29,0.53}{#1}}
\newcommand{\DecValTok}[1]{\textcolor[rgb]{0.00,0.00,0.81}{#1}}
\newcommand{\DocumentationTok}[1]{\textcolor[rgb]{0.56,0.35,0.01}{\textbf{\textit{#1}}}}
\newcommand{\ErrorTok}[1]{\textcolor[rgb]{0.64,0.00,0.00}{\textbf{#1}}}
\newcommand{\ExtensionTok}[1]{#1}
\newcommand{\FloatTok}[1]{\textcolor[rgb]{0.00,0.00,0.81}{#1}}
\newcommand{\FunctionTok}[1]{\textcolor[rgb]{0.13,0.29,0.53}{\textbf{#1}}}
\newcommand{\ImportTok}[1]{#1}
\newcommand{\InformationTok}[1]{\textcolor[rgb]{0.56,0.35,0.01}{\textbf{\textit{#1}}}}
\newcommand{\KeywordTok}[1]{\textcolor[rgb]{0.13,0.29,0.53}{\textbf{#1}}}
\newcommand{\NormalTok}[1]{#1}
\newcommand{\OperatorTok}[1]{\textcolor[rgb]{0.81,0.36,0.00}{\textbf{#1}}}
\newcommand{\OtherTok}[1]{\textcolor[rgb]{0.56,0.35,0.01}{#1}}
\newcommand{\PreprocessorTok}[1]{\textcolor[rgb]{0.56,0.35,0.01}{\textit{#1}}}
\newcommand{\RegionMarkerTok}[1]{#1}
\newcommand{\SpecialCharTok}[1]{\textcolor[rgb]{0.81,0.36,0.00}{\textbf{#1}}}
\newcommand{\SpecialStringTok}[1]{\textcolor[rgb]{0.31,0.60,0.02}{#1}}
\newcommand{\StringTok}[1]{\textcolor[rgb]{0.31,0.60,0.02}{#1}}
\newcommand{\VariableTok}[1]{\textcolor[rgb]{0.00,0.00,0.00}{#1}}
\newcommand{\VerbatimStringTok}[1]{\textcolor[rgb]{0.31,0.60,0.02}{#1}}
\newcommand{\WarningTok}[1]{\textcolor[rgb]{0.56,0.35,0.01}{\textbf{\textit{#1}}}}
\usepackage{graphicx}
\makeatletter
\newsavebox\pandoc@box
\newcommand*\pandocbounded[1]{% scales image to fit in text height/width
  \sbox\pandoc@box{#1}%
  \Gscale@div\@tempa{\textheight}{\dimexpr\ht\pandoc@box+\dp\pandoc@box\relax}%
  \Gscale@div\@tempb{\linewidth}{\wd\pandoc@box}%
  \ifdim\@tempb\p@<\@tempa\p@\let\@tempa\@tempb\fi% select the smaller of both
  \ifdim\@tempa\p@<\p@\scalebox{\@tempa}{\usebox\pandoc@box}%
  \else\usebox{\pandoc@box}%
  \fi%
}
% Set default figure placement to htbp
\def\fps@figure{htbp}
\makeatother
\setlength{\emergencystretch}{3em} % prevent overfull lines
\providecommand{\tightlist}{%
  \setlength{\itemsep}{0pt}\setlength{\parskip}{0pt}}
\usepackage{bookmark}
\IfFileExists{xurl.sty}{\usepackage{xurl}}{} % add URL line breaks if available
\urlstyle{same}
\hypersetup{
  pdftitle={DBSCAN as an ML Tool},
  hidelinks,
  pdfcreator={LaTeX via pandoc}}

\title{DBSCAN as an ML Tool}
\author{}
\date{\vspace{-2.5em}2026-04-17}

\begin{document}
\maketitle

\subsection{Introduction}\label{introduction}

DBSCAN (Density-Based Spatial Clustering of Applications with Noise)
groups points that are densely packed together and flags isolated points
as noise. Unlike \(k\)-means it does not require specifying the number
of clusters, handles non-convex cluster shapes, and produces an
outlier-detection signal as a side effect --- every point flagged as
noise is a candidate anomaly. For machine-learning workflows where
natural cluster discovery and outlier flagging are both useful, DBSCAN
is the right starting point.

\subsection{Prerequisites}\label{prerequisites}

A working understanding of distance metrics, density-based reasoning,
and the limitations of partitioning algorithms (\(k\)-means) for
non-convex cluster shapes.

\subsection{Theory}\label{theory}

DBSCAN has two hyperparameters: \texttt{eps} (the neighbourhood radius)
and \texttt{minPts} (the minimum number of points required for a dense
region). Three categories of point emerge from the algorithm:

\begin{itemize}
\tightlist
\item
  \textbf{Core point}: at least \texttt{minPts} neighbours within
  \texttt{eps}.
\item
  \textbf{Border point}: within \texttt{eps} of a core point but not
  itself a core point.
\item
  \textbf{Noise point}: neither core nor border.
\end{itemize}

Clusters are formed as maximally connected sets of core points (with
their associated border points). Choosing \texttt{eps} is the practical
challenge; the standard heuristic is the \(k\)-distance plot, which
shows the distance to the \(k\)-th nearest neighbour for each point
sorted ascending --- a knee in the curve marks a sensible \texttt{eps}.

\subsection{Assumptions}\label{assumptions}

Cluster density is roughly uniform; the chosen \texttt{eps} matches the
natural scale of the data; Euclidean distance is meaningful (or another
distance has been substituted).

\subsection{R Implementation}\label{r-implementation}

\begin{Shaded}
\begin{Highlighting}[]
\FunctionTok{library}\NormalTok{(dbscan)}

\FunctionTok{set.seed}\NormalTok{(}\DecValTok{2026}\NormalTok{)}
\NormalTok{n }\OtherTok{\textless{}{-}} \DecValTok{200}
\NormalTok{x1 }\OtherTok{\textless{}{-}} \FunctionTok{c}\NormalTok{(}\FunctionTok{rnorm}\NormalTok{(n}\SpecialCharTok{/}\DecValTok{2}\NormalTok{, }\SpecialCharTok{{-}}\DecValTok{2}\NormalTok{, }\FloatTok{0.3}\NormalTok{), }\FunctionTok{rnorm}\NormalTok{(n}\SpecialCharTok{/}\DecValTok{2}\NormalTok{, }\DecValTok{2}\NormalTok{, }\FloatTok{0.3}\NormalTok{))}
\NormalTok{x2 }\OtherTok{\textless{}{-}} \FunctionTok{c}\NormalTok{(}\FunctionTok{rnorm}\NormalTok{(n}\SpecialCharTok{/}\DecValTok{2}\NormalTok{, }\SpecialCharTok{{-}}\DecValTok{2}\NormalTok{, }\FloatTok{0.3}\NormalTok{), }\FunctionTok{rnorm}\NormalTok{(n}\SpecialCharTok{/}\DecValTok{2}\NormalTok{, }\DecValTok{2}\NormalTok{, }\FloatTok{0.3}\NormalTok{))}
\NormalTok{outliers\_x1 }\OtherTok{\textless{}{-}} \FunctionTok{runif}\NormalTok{(}\DecValTok{10}\NormalTok{, }\SpecialCharTok{{-}}\DecValTok{4}\NormalTok{, }\DecValTok{4}\NormalTok{)}
\NormalTok{outliers\_x2 }\OtherTok{\textless{}{-}} \FunctionTok{runif}\NormalTok{(}\DecValTok{10}\NormalTok{, }\SpecialCharTok{{-}}\DecValTok{4}\NormalTok{, }\DecValTok{4}\NormalTok{)}
\NormalTok{df }\OtherTok{\textless{}{-}} \FunctionTok{data.frame}\NormalTok{(}\AttributeTok{x1 =} \FunctionTok{c}\NormalTok{(x1, outliers\_x1),}
                 \AttributeTok{x2 =} \FunctionTok{c}\NormalTok{(x2, outliers\_x2))}

\CommentTok{\# Choose eps via the k{-}distance plot knee}
\FunctionTok{kNNdistplot}\NormalTok{(df, }\AttributeTok{k =} \DecValTok{5}\NormalTok{); }\FunctionTok{abline}\NormalTok{(}\AttributeTok{h =} \FloatTok{0.5}\NormalTok{, }\AttributeTok{lty =} \DecValTok{2}\NormalTok{, }\AttributeTok{col =} \StringTok{"red"}\NormalTok{)}

\NormalTok{db }\OtherTok{\textless{}{-}} \FunctionTok{dbscan}\NormalTok{(df, }\AttributeTok{eps =} \FloatTok{0.5}\NormalTok{, }\AttributeTok{minPts =} \DecValTok{5}\NormalTok{)}
\FunctionTok{table}\NormalTok{(}\AttributeTok{cluster =}\NormalTok{ db}\SpecialCharTok{$}\NormalTok{cluster)}

\FunctionTok{plot}\NormalTok{(df, }\AttributeTok{col =}\NormalTok{ db}\SpecialCharTok{$}\NormalTok{cluster }\SpecialCharTok{+} \DecValTok{1}\NormalTok{, }\AttributeTok{pch =} \DecValTok{16}\NormalTok{, }\AttributeTok{asp =} \DecValTok{1}\NormalTok{,}
     \AttributeTok{main =} \StringTok{"DBSCAN clustering (cluster 0 = noise)"}\NormalTok{,}
     \AttributeTok{xlab =} \StringTok{"x1"}\NormalTok{, }\AttributeTok{ylab =} \StringTok{"x2"}\NormalTok{)}
\end{Highlighting}
\end{Shaded}

\subsection{Output \& Results}\label{output-results}

\texttt{dbscan()} returns a vector of cluster labels with \texttt{0}
indicating noise. The \(k\)-distance plot guides \texttt{eps} selection;
the cluster scatter shows densely-packed groups separated by noise
points.

\subsection{Interpretation}\label{interpretation}

A reporting sentence: ``DBSCAN with \texttt{eps\ =\ 0.5} and
\texttt{minPts\ =\ 5} recovered the two natural clusters injected in the
simulation (96 of 100 cluster-1 points and 98 of 100 cluster-2 points
correctly assigned) and flagged 10 of the 10 injected outliers as noise
(cluster 0); the remaining noise calls were 6 boundary points where the
local density is genuinely lower.'' Always state \texttt{eps},
\texttt{minPts}, and the noise count.

\subsection{Practical Tips}\label{practical-tips}

\begin{itemize}
\tightlist
\item
  Set \texttt{eps} via the \(k\)-distance plot's knee; it is the most
  sensitive hyperparameter and a poor choice fails the algorithm
  completely.
\item
  `minPts = 2 \times d\$ (dimensions) is a sensible starting heuristic;
  inspect cluster counts and adjust if they are unreasonably few or
  many.
\item
  For varying-density clusters, switch to HDBSCAN
  (\texttt{dbscan::hdbscan}); it adapts the density threshold per region
  and is more robust.
\item
  Use the noise label (\texttt{cluster\ ==\ 0}) as a built-in outlier
  detector; for anomaly-detection tasks, this is often the simplest and
  most interpretable approach.
\item
  Scale features before clustering (\texttt{scale()} or
  \texttt{recipes::step\_normalize()}); DBSCAN's Euclidean distance is
  sensitive to feature scales.
\item
  For high-dimensional data, distance concentration weakens DBSCAN;
  reduce dimensionality with PCA or UMAP first.
\end{itemize}

\subsection{R Packages Used}\label{r-packages-used}

\texttt{dbscan} for \texttt{dbscan()}, \texttt{hdbscan()},
\texttt{optics()}, and the \(k\)-distance heuristic;
\texttt{fpc::dbscan()} for an alternative implementation;
\texttt{cluster} for related clustering algorithms;
\texttt{factoextra::fviz\_cluster()} for ggplot-style visualisation of
DBSCAN output.

\subsection{For Reviewers}\label{for-reviewers}

\textbf{What to look for in a paper using this method.}

\begin{itemize}
\tightlist
\item
  \textbf{Common misapplications.}
\item
  \textbf{Diagnostics that should be reported but often aren't.}
\item
  \textbf{Red flags in tables and figures.}
\item
  \textbf{What to verify.}
\item
  \textbf{What an adequate Methods paragraph must contain.}
\end{itemize}

\subsection{See also --- labs in this
chapter}\label{see-also-labs-in-this-chapter}

\begin{itemize}
\tightlist
\item
  \href{labs/lab-cross-validation-nested-cv-bootstrap-632.Rmd}{Cross-validation,
  nested CV, bootstrap .632+}
\item
  \href{labs/lab-regularisation-ridge-lasso-elastic-net.Rmd}{Regularisation:
  ridge, lasso, elastic net}
\item
  \href{labs/lab-trees-random-forests-gradient-boosting.Rmd}{Trees,
  random forests, gradient boosting}
\item
  \href{labs/lab-interpretability-and-shap.Rmd}{Interpretability and
  SHAP}
\item
  \href{labs/lab-tabular-neural-networks-with-torch.Rmd}{Tabular neural
  networks with torch}
\item
  \href{labs/lab-imaging-and-sequence-models-intro.Rmd}{Imaging and
  sequence models intro}
\item
  \href{labs/lab-tidymodels-pipelines.Rmd}{tidymodels pipelines}
\item
  \href{labs/lab-fdr-knockoffs-replication-crisis.Rmd}{FDR, knockoffs,
  replication crisis}
\end{itemize}

\end{document}
